% Challenge 1: state the goal and proof tasks; keep the expansion in the notes.
\begin{frame}{RSP Challenge 1: Preparation Errors}
  \centering
  \setbeamercolor{itemize item}{fg=CVQCTeal}
  \setbeamercolor{itemize subitem}{fg=CVQCTeal}
  \begin{tikzpicture}[
      >={Latex[length=2.5mm]},
      protocol/.style={draw=CVQCInk,fill=CVQCPale,thick,
        rounded corners=2mm,minimum width=1.7cm,minimum height=0.80cm,
        inner sep=3pt,font=\large,text=CVQCInk},
      output/.style={inner sep=0pt,font=\large,text=CVQCInk},
      flow/.style={->,very thick,draw=CVQCTeal}
    ]
    \path[use as bounding box] (-6.6,-3.15) rectangle (6.6,3.05);

    \foreach \x in {-5.1,-2.5,0.1,5.1} {
      \node[protocol] at (\x,2.10) {$\pi_{\mathrm{RSP}}$};
      \draw[flow] (\x,1.70) -- (\x,1.22);
    }
    \node[font=\large,text=CVQCMuted] at (2.7,2.10) {$\cdots$};
    \node[font=\large,text=CVQCMuted] at (2.7,0.88) {$\cdots$};

    \node[output] at (-5.1,0.88) {$\ket{+_{\theta_1}}$};
    \node[output] at (-2.5,0.88) {$\ket{+_{\theta_2}}$};
    \node[output] at (0.1,0.88) {$\ket{+_{\theta_3}}$};
    \node[output] at (5.1,0.88) {$\ket{+_{\theta_K}}$};

    \node[anchor=north west,inner sep=0pt,text=CVQCInk] at (-6.1,0.30) {%
      \begin{minipage}{12.2cm}
        \raggedright
        \begin{itemize}
          \setlength{\itemsep}{0.55em}
          \item<1-> Each RSP output incurs $1/\mathrm{poly}$ error,
            possibly coherent across copies.
          \item<1-> \textbf{Goal:} distillable inputs---most outputs are good.
          \item<2-> \textbf{Analysis needs to:}
            \begin{itemize}
              \setlength{\itemsep}{0.20em}
              \item<2-> Upper-bound high-weight attack components.
              \item<2-> Recover factorizability.
            \end{itemize}
        \end{itemize}
      \end{minipage}%
    };
  \end{tikzpicture}

  \note<1>{
    \footnotesize
    The row shows intended preparation states. It does not assert independent
    calls or a tensor-product output against a malicious prover. The invocations
    run sequentially, interleaved with hidden tests.
    The first bullet is informal shorthand for the inverse-polynomial
    approximation estimate, with the conjugation caveat on the next slide.
    Coherent errors are not independently sampled faulty-qubit events.
    Directly replacing the real preparations by approximate ideal states would
    carry an inverse-polynomial error into the final result. Instead, write each
    real preparation map as its exact ideal counterpart plus a difference map.
    Both maps include the prover's memory and previously prepared registers,
    so this allows correlated attacks. The error map need not be positive or
    a physical operation. Expand chronologically and stop a term when it
    reaches d error factors.
    Source: the paper's Classical Remote State Preparation section,
    ideal channel construction and adaptive error expansion.
  }

  \note<2>{
    \scriptsize
    Many-error terms: the first encrypted response precedes revelation of
    whether the call is a test or a preparation. Encryption hiding, the
    self-testing estimate, and a sequential concentration argument bound the
    accepted contribution of terms reaching d errors, against efficient
    distinguishers. Polynomially many tests suffice.
    Few-error terms: if their error factors occur on calls S, only the MBQC
    copies receiving those inputs can acquire additional plaintext-level
    faults from RSP. Their number is at most the size of S.
    These error maps may depend on the private preparation descriptions,
    so earlier factorizability cannot be assumed. Fourier expansion and
    the MBQC identities establish it while keeping the remaining attack
    on all copies and the prover's memory.
    Cut-and-choose bounds additional MBQC faults. Distillation removes
    plaintext-level faults, and authenticated evaluation controls further
    deviations. Combining signed components and their cross terms costs
    an exponential factor. The parameter choice makes the earlier
    information-theoretic error small enough to absorb that factor.
    This is an expansion used in the proof, not a probabilistic claim that
    the real source chooses a small set of faulty physical qubits.
    Source: source summary and composition theorems in the paper.
  }

\end{frame}
